\scanpage{119}
\mspage{- 90 -}
n\textsuperscript{o} 371

% L1–2
tout $i$. Montrons que l'on a $(u(x))^2=0$ pour tout $x\in M$ ; on peut en effet écrire $x=\sum_{i=1}^{n}x_i$, où $x_i\in M_i$, donc $(u(x))^2=\sum_{i=1}^{n}(u(x_i))^2+\sum_{i<j}(u(x_i)u(x_j)+u(x_j)u(x_i))$.
% L3–6
Or, comme $\handwritten{\psi}_i(x_i)$ est un élément de degré 1 dans $\handwritten{\bigwedge}(M_i)$, on a, pour $i\ne j$, $h_i(\handwritten{\psi}_i(x_i))h_j(\handwritten{\psi}_j(x_j))=-h_j(\handwritten{\psi}_j(x_j))h_i(\handwritten{\psi}_i(x_i))$ (\S~3, n\textsuperscript{o}3, prop.2), d'où $u(x_i)u(x_j)+u(x_j)u(x_i)=0$ ; par ailleurs $(u(x_i))^2=h_i((\handwritten{\psi}_i(x_i))^2)=0$ par définition de $\handwritten{\bigwedge}(M_i)$. En vertu de la prop.3 du n\textsuperscript{o}2, il existe donc un homomorphisme unique d'algèbres graduées
% L7–8
$h:\handwritten{\bigwedge}(M)\longrightarrow {}^{\mathrm g}\!\bigotimes_i(\handwritten{\bigwedge}(M_i))$ tel que $u=h\circ\handwritten{\psi}$ (où $\handwritten{\psi}:M\longrightarrow\handwritten{\bigwedge}(M)$ est l'injection canonique). Pour $x_i\in M_i$, on a
% L9–12
$h(g(u_i(x_i)))=h(g(h_i(\handwritten{\psi}_i(x_i))))=h(f_i(\handwritten{\psi}_i(x_i)))=h(\handwritten{\psi}(j_i(x_i))=u_i(x_i)$, et comme les $u_i(x_i)$ engendrent l'algèbre $ {}^{\mathrm g}\!\bigotimes_i(\handwritten{\bigwedge}(M_i))$, $h\circ g$ est l'identité ; de même $g(h(\handwritten{\psi}(j_i(x_i))))=g(u_i(x_i))=g(h_i(\handwritten{\psi}_i(x_i)))=f_i(\handwritten{\psi}_i(x_i))=\handwritten{\psi}(j_i(x_i))$, et les éléments $\handwritten{\psi}(j_i(x_i))$ engendrent l'algèbre $\handwritten{\bigwedge}(M)$, donc $g\circ h$ est l'identité.

% L13–16
COROLLAIRE 1.- \uline{Sous les hypothèses de la prop.7}, \uline{pour tout entier} $m\geq 0$, \uline{la puissance extérieure m-ème} $\overset{m}{\handwritten{\bigwedge}}\,M$ \uline{s'identifie canoniquement par} $g$ \uline{à la somme directe des modules} $\bigotimes_{1\leq i\leq n}(\overset{m_i}{\handwritten{\bigwedge}}\,M_i)$, \uline{où} $(m_i)_{1\leq i\leq n}$ \uline{parcourt l'ensemble des suites d'entiers} $\geq 0$ \uline{telles que} $\sum_{i=1}^{n}m_i=m$.

% L17–19
Cet isomorphisme s'explicite de la façon suivante : considérons $m$ éléments $x_j=\sum_{i=1}^{n}x_{ji}$ $(1\leq j\leq m)$ de $M$, avec $x_{ji}\in M_i$ pour $1\leq i\leq n$, et pour toute application $\handwritten{\sigma}$ de $\unsure{[1,m]|[\![1,m]\!]}$ dans $[1,n]$, formons le m-vecteur
% L20
\begingroup\csname tagsleft@true\endcsname
\[
 x_{1,\handwritten{\sigma}(1)}\handwritten{\wedge} x_{2,\handwritten{\sigma}(2)}\handwritten{\wedge}\ldots\handwritten{\wedge} x_{m,\handwritten{\sigma}(m)}\quad .\tag*{(7)}
\]
\endgroup
% L21–24
La donnée de l'application $\handwritten{\sigma}$ équivaut à la donnée, pour chaque $i\in[1,n]$ d'une partie $H_i=\handwritten{\sigma}^{-1}(i)$ de $[1,m]$, celles des $H_i$ non vides formant une \uline{partition} de $[1,m]$ ; si $H_i$ est vide, posons $x_{H_i,\,i}=1$, et dans le cas contraire, $x_{H_i,\,i}=x_{j_1,\,i}\handwritten{\wedge} x_{j_2,\,i}\handwritten{\wedge}\ldots\handwritten{\wedge} x_{j_{r_i},\,i}$, où $(j_k)_{1\leq k\leq r_i}$ est la suite
% archive sheet number: 118
